% !TEX program = xelatex
\documentclass[11pt,a4paper]{article}

% ── Packages ──────────────────────────────────────────────────────────────────
\usepackage[margin=2.5cm]{geometry}
\usepackage{booktabs}
\usepackage{graphicx}
\usepackage[colorlinks=true,linkcolor=blue!70!black,citecolor=blue!70!black,urlcolor=blue!70!black]{hyperref}
\usepackage{amsmath}
\usepackage[round]{natbib}
\usepackage{xcolor}
\usepackage{multirow}
\usepackage{float}
\usepackage{enumitem}
\usepackage{array}
\usepackage{tabularx}
\usepackage{xeCJK}
\setCJKmainfont{STFangsong}
\usepackage{algorithm}
\usepackage{algorithmic}
\usepackage{caption}
\usepackage{subcaption}
\usepackage{fancyhdr}
\usepackage{setspace}
\newcommand{\yen}{¥}

% ── Page style ────────────────────────────────────────────────────────────────
\pagestyle{fancy}
\fancyhf{}
\fancyhead[L]{\small\thepage}
\fancyhead[R]{\small AtomGradient}
\renewcommand{\headrulewidth}{0.4pt}

% ── Custom commands ───────────────────────────────────────────────────────────
\newcommand{\echostream}{\textsc{Prism}}
\newcommand{\iir}{\textsc{IIR}}
\newcommand{\appname}[1]{\textsf{#1}}
\newcommand{\tiernum}[1]{\textbf{Tier~#1}}

% ── Title ─────────────────────────────────────────────────────────────────────
\title{%
  \textbf{Prism: Cross-Domain Personal Data Integration on \\
  Consumer Hardware Produces Emergent Insights}%
}

\author{%
  AtomGradient\\
  \texttt{\url{https://github.com/AtomGradient/Prism}}\\
  \texttt{\url{https://atomgradient.github.io/Prism}}
}

\date{March 2026}

\begin{document}

\maketitle

% ══════════════════════════════════════════════════════════════════════════════
% ABSTRACT
% ══════════════════════════════════════════════════════════════════════════════
\begin{abstract}
Current cloud-based AI assistants operate under a fundamental structural
contradiction: to be genuinely useful, a model must understand its user's life
context, yet the private data needed for such understanding is precisely what
users are unwilling to upload to remote servers. We present \echostream{}, a
three-tier on-device architecture that resolves this privacy--utility paradox by
integrating data from four vertical personal-data applications---finance, diet,
mood, and reading---entirely on consumer Apple Silicon hardware. Through an
ablation study on 10 synthetic users over 90 simulated days, we demonstrate
that full cross-domain data integration achieves an average Insight Increment
Ratio (\iir{}) of $1.48\times$ over single-domain baselines, with the
cross-domain scoring dimension driving nearly all of the gain. A model-scale
study across four Qwen3.5 checkpoints (0.8B to 35B-A3B) reveals diminishing
returns: the 2B$\to$9B jump yields $+14.9$ points while 9B$\to$35B yields only
$+5.6$. Our federation protocol transmits only data summaries, achieving a
$125.5\times$ compression ratio with zero raw-data leakage. All experiments run
on commodity hardware (M2~Ultra, M1~Max, M2~Pro), establishing a practical
blueprint for privacy-preserving personal AI.
\end{abstract}

\vspace{0.5em}
\noindent\textbf{Keywords:} personal AI, cross-domain data integration,
on-device inference, federated data protocol, Apple Silicon, privacy-preserving AI

% ══════════════════════════════════════════════════════════════════════════════
\section{Introduction}
\label{sec:intro}
% ══════════════════════════════════════════════════════════════════════════════

\subsection{The Privacy--Utility Paradox}

The promise of a truly personal AI assistant rests on a circular dependency
that current cloud-centric architectures cannot resolve: \emph{to make AI
genuinely useful, it must understand its user; to understand a user, it needs
their most private data; yet the more private the data, the less willing users
are to upload it to a third-party server.}

Today's flagship models---GPT-4o, Claude, Gemini---are powerful
general-purpose reasoners, but they remain \emph{knowledgeable strangers}.
They can draft emails, generate images, and answer trivia, yet they do not
know whether their user slept well last night, is under financial stress, or
has been skipping meals for a week. The data that would enable such
understanding---journals, medical records, financial transactions, dietary
logs, emotional states---is exactly what users refuse to entrust to remote
servers.

This is not merely a technical obstacle; it is a \emph{structural} one. Cloud
AI providers derive revenue from API calls (per-token billing); data remaining
on-device eliminates their primary business incentive. Their competitive moat
is built on data accumulation; local-only data erodes that moat. Even
platform-native players face barriers: Apple's app ecosystem fragments
third-party data into isolated silos, and Google's advertising model creates a
rational distrust of deeper data access.

\subsection{A Paradigm Inversion}

We observe that the dominant paradigm can be characterized as \textbf{large
model + small data}: a cloud-hosted model with hundreds of billions of
parameters operates on a thin sliver of user context squeezed into a few
thousand tokens of conversation history. \echostream{} inverts this paradigm
to \textbf{medium model + rich data}: a locally hosted model of modest scale
(3--35 billion parameters) operates on the user's entire longitudinal life
data across multiple domains.

The core insight is that intelligence does not arise solely from parameter
count; it also arises from \emph{data depth and authenticity}. A 9-billion
parameter model that has access to 90 days of a user's financial transactions,
dietary patterns, emotional states, and reading habits can produce insights
that a 1-trillion parameter model operating on a single conversation turn
cannot---because the latter simply does not have the data.

\subsection{Contributions}

This paper makes the following contributions:

\begin{enumerate}[leftmargin=*,itemsep=2pt]
  \item \textbf{Cross-domain insight emergence.} We demonstrate through a
    controlled ablation study (10 users $\times$ 8 data configurations) that
    integrating four personal-data domains produces an average \iir{} of
    $1.48\times$ over single-domain analysis, with the cross-domain scoring
    dimension accounting for nearly all of the gain (Section~\ref{sec:results:ablation}).

  \item \textbf{Model-scale analysis.} We establish a model-size vs.\
    insight-quality curve on commodity hardware, showing that the practical
    sweet spot lies at the 9B parameter scale, with diminishing returns beyond
    that point (Section~\ref{sec:results:scale}).

  \item \textbf{Privacy-preserving federation.} We implement and evaluate a
    LAN-based federation protocol that achieves $125.5\times$ data compression
    with zero raw-data leakage, demonstrating that multi-device personal AI can
    operate without any data leaving the user's physical premises
    (Section~\ref{sec:results:federation}).

  \item \textbf{Practical system design.} We provide a fully reproducible
    three-tier architecture running entirely on consumer Apple Silicon hardware,
    with all code and synthetic data publicly available.\footnote{%
      \url{https://github.com/AtomGradient/Prism}}
\end{enumerate}


% ══════════════════════════════════════════════════════════════════════════════
\section{Related Work}
\label{sec:related}
% ══════════════════════════════════════════════════════════════════════════════

\subsection{Personalized Large Language Models}

The challenge of adapting LLMs to individual users has received growing
attention. \citet{survey_personalized_llm_2025} provide a comprehensive survey
identifying three axes of personalization: user profile injection, few-shot
behavioral conditioning, and parametric adaptation. Most approaches assume
access to a relatively small and static user profile.

\citet{plum2024} propose PLUM, a system that learns to ``remember'' user
conversations by compressing dialogue history into persistent user embeddings.
While effective for conversational style adaptation, PLUM does not address
cross-domain data integration or the privacy implications of storing user
embeddings on cloud infrastructure.

PersonalLLM \citep{personalllm2025} introduces a benchmark and training
methodology for personalized generation, demonstrating that even modest
amounts of user-specific data can significantly improve output alignment.
However, the evaluation is confined to text generation quality within a single
domain and does not consider the compounding effects of multi-domain data
fusion.

\subsection{Continual and Lifelong Learning}

The temporal dynamics of user preferences---what the literature terms
\emph{temporal drift}---pose a distinct challenge. \citet{continual_llm_survey_2024}
survey continual learning methods for LLMs, highlighting catastrophic
forgetting as the primary obstacle to lifelong adaptation. Current approaches
(elastic weight consolidation, progressive neural networks, replay buffers)
focus on model-side solutions and do not consider the data architecture that
feeds the model.

LIBER \citep{liber2024} addresses lifelong user behavior modeling by proposing
a framework that compresses long behavioral sequences into structured
representations. LIBER's insight---that raw behavioral sequences are too long
for direct LLM consumption and require intermediate structuring---directly
informs our federation protocol design. However, LIBER operates in a
cloud-based recommendation setting and does not address on-device deployment
or cross-domain integration.

\subsection{Federated Learning and On-Device Inference}

Federated learning \citep{mcmahan2017fedavg} was designed to train models
across distributed devices without centralizing raw data. However, classical
federated learning focuses on \emph{model training} (gradient aggregation),
whereas our setting requires \emph{data-level integration} for inference-time
reasoning. Our federation protocol is closer to federated analytics
\citep{ramage2020federated_analytics} but operates at the semantic summary
level rather than the statistical aggregate level.

On-device LLM inference has become practical with the advent of Apple Silicon's
unified memory architecture. The MLX framework \citep{mlx2023} and
llama.cpp \citep{llamacpp2023} have demonstrated that models up to 70B
parameters (with quantization) can run at interactive speeds on consumer
hardware. Our work builds on this capability but focuses on \emph{what to do}
with on-device inference once it is technically feasible, rather than on the
inference engine itself.


% ══════════════════════════════════════════════════════════════════════════════
\section{System Design}
\label{sec:system}
% ══════════════════════════════════════════════════════════════════════════════

\subsection{Design Principles}

\echostream{} is built on three design principles:

\begin{enumerate}[leftmargin=*,itemsep=2pt]
  \item \textbf{Privacy by architecture, not by policy.} Raw data never
    leaves the device that generated it. This is enforced structurally, not
    by terms of service.
  \item \textbf{Data collection as a natural byproduct.} Users interact with
    genuinely useful vertical applications; data accumulation is a side effect
    of normal usage, not a separate act of ``feeding the AI.''
  \item \textbf{Heterogeneous device awareness.} The system adapts its
    inference strategy to the computational tier of the device, from sub-1B
    models on mobile to 35B MoE models on a home server.
\end{enumerate}

\subsection{Three-Tier Device Topology}

The system organizes user devices into three computational tiers:

\paragraph{Tier 0 --- Mobile Devices (Data Collection + Lightweight Inference).}
Smartphones and tablets serve as the primary data collection points. In our
experimental setup, an iPhone 15 Pro Max (8\,GB) runs \appname{Mealens}
(dietary tracking via photo recognition) and \appname{Ururu} (mood journaling
and emotional state tracking), while an iPad Air M3 (8\,GB) runs
\appname{Narrus} (reading analysis) and \appname{Dailyn} (financial
accounting). These devices can run sub-2B models for immediate local queries
but delegate complex cross-domain reasoning to higher tiers.

\paragraph{Tier 1 --- Mid-Tier Local Inference (Daily Queries).}
Machines with 32\,GB unified memory (Apple M1~Max, M2~Pro) run quantized 9B
models, handling routine queries that require moderate reasoning over a single
user's data. In our experiments, an M1~Max 32\,GB and an M2~Pro 32\,GB serve
this role, achieving 21.9 and 12.7 tokens per second respectively on
Qwen3.5-9B-Q8.

\paragraph{Tier 2 --- Home Server (Panoramic Inference + Model Evolution).}
A machine with large unified memory (Apple M2~Ultra 192\,GB) runs
full-precision or lightly quantized 35B-class models. This tier performs
cross-domain ``panoramic'' reasoning by federating data summaries from all
Tier~0 devices. In idle periods, it can execute LoRA incremental fine-tuning
on Tier~1 models using accumulated data. In our experiments, the M2~Ultra runs
Qwen3.5-35B-A3B-Q8 at 49.9 tokens per second.

\subsection{Vertical Application Suite}

Four applications collect data from complementary life domains:

\begin{table}[H]
\centering
\caption{Vertical application suite and the data dimensions they capture.}
\label{tab:apps}
\begin{tabular}{@{}llll@{}}
\toprule
\textbf{App} & \textbf{Domain} & \textbf{Data Dimension} & \textbf{Example Signals} \\
\midrule
\appname{Dailyn}  & Finance    & Behavioral layer & Spending categories, amounts, timing \\
\appname{Mealens} & Diet       & Physical layer   & Meal photos, nutritional estimates \\
\appname{Ururu}   & Mood       & Emotional layer  & Mood scores, journal entries, triggers \\
\appname{Narrus}  & Reading    & Cognitive layer  & Articles read, topics, time spent \\
\bottomrule
\end{tabular}
\end{table}

The critical design decision is that these are \emph{standalone useful tools},
not data-collection wrappers. A user adopts \appname{Dailyn} because it is a
good accounting app; the fact that it contributes to a cross-domain personal
model is invisible and incidental. This resolves the cold-start problem: data
accumulates as a natural byproduct of tool usage.

\subsection{Federation Protocol}
\label{sec:system:federation}

When a Tier~2 device initiates a panoramic query, the federation protocol
executes the following steps:

\begin{algorithm}[H]
\caption{LAN Federation Protocol for Panoramic Query}
\label{alg:federation}
\begin{algorithmic}[1]
\REQUIRE Query $q$, user ID $u$, time window $\Delta t$, registered data nodes $\mathcal{N}$
\ENSURE Panoramic insight $\mathcal{I}$
\STATE $\mathcal{S} \leftarrow \{\}$ \COMMENT{Collected summaries}
\FOR{each node $n \in \mathcal{N}$ \textbf{in parallel}}
  \STATE $s_n \leftarrow$ \textsc{RequestSummary}($n$, $u$, $\Delta t$) \COMMENT{LAN-local HTTP}
  \STATE $\mathcal{S} \leftarrow \mathcal{S} \cup \{s_n\}$
\ENDFOR
\STATE $\text{prompt} \leftarrow$ \textsc{BuildPanoramicPrompt}($q$, $\mathcal{S}$)
\STATE $\mathcal{I} \leftarrow$ \textsc{LLMInference}($\text{prompt}$) \COMMENT{Tier 2 model}
\STATE \textsc{AuditLog}($|\text{raw}|$, $|\text{transmitted}|$, $\text{nodes}$) \COMMENT{Privacy audit}
\RETURN $\mathcal{I}$
\end{algorithmic}
\end{algorithm}

Each data node exposes a standardized JSON Schema interface that returns
\emph{summaries}, never raw records. For example, \appname{Dailyn}'s summary
endpoint returns aggregated spending by category and trend indicators, not
individual transaction records. The privacy audit trail records the size of
raw data on each device versus the bytes actually transmitted, enabling
post-hoc verification that no data leakage occurred.

In our experiments, the protocol achieves a $125.5\times$ compression ratio:
108,850 bytes of raw data across all apps compress to 867 bytes of transmitted
summaries. The average federation latency (network round-trip to all data
nodes) is 469.5\,ms, negligible compared to the LLM inference time of 40.3\,s.


% ══════════════════════════════════════════════════════════════════════════════
\section{Experimental Setup}
\label{sec:setup}
% ══════════════════════════════════════════════════════════════════════════════

\subsection{Synthetic User Population}

We generate synthetic data for 10 users spanning diverse socioeconomic
profiles, each with 90 days of simulated life data across all four
applications. Each user has a pre-designed ``life event'' injected at a
specific day to simulate temporal drift---the phenomenon where a significant
event causes correlated changes across multiple life domains simultaneously.

\begin{table}[H]
\centering
\caption{Synthetic user profiles with injected life events for temporal drift simulation.}
\label{tab:users}
\small
\begin{tabular}{@{}clccp{5.5cm}@{}}
\toprule
\textbf{ID} & \textbf{Profile} & \textbf{Income} & \textbf{Age} & \textbf{Injected Event} \\
\midrule
user\_01 & Factory worker     & \yen 3,500/mo  & 22 & Day 40: factory layoff, income interrupted \\
user\_02 & Delivery rider     & \yen 6,000/mo  & 28 & Day 46: severe traffic accident $\to$ hospitalization $\to$ debt spiral \\
user\_03 & Primary teacher    & \yen 5,500/mo  & 31 & Day 35: pregnancy confirmed, spending/diet shift \\
user\_04 & Grad student       & \yen 2,000/mo  & 26 & Day 45: thesis rejected by advisor \\
user\_05 & E-commerce ops     & \yen 12,000/mo & 27 & Day 30: promoted to team lead \\
user\_06 & Freelance illustrator & \yen 8--25K/mo & 29 & Day 60: large client commission (\yen 50K) \\
user\_07 & Hospital resident  & \yen 25,000/mo & 33 & Day 21: medication error $\to$ patient ICU $\to$ suspension $\to$ breakup \\
user\_08 & Tech P7 engineer   & \yen 50,000/mo & 32 & Day 55: laid off (N+1) $\to$ divorce $\to$ depression spiral \\
user\_09 & Business owner & \yen 150,000/mo & 45 & Day 38: \yen 800K bad debt $\to$ cash flow crisis $\to$ hospitalization \\
user\_10 & Fund manager       & \yen 830,000/mo & 38 & Day 25: 4\% single-day fund drawdown \\
\bottomrule
\end{tabular}
\end{table}

The user population is deliberately diverse in income level (from
\yen 2,000/month graduate student to \yen 830,000/month fund manager), life
stage (22-year-old factory worker to 45-year-old business owner), and event
type (sudden income loss, health changes, career shifts). This diversity tests
whether cross-domain insight emergence is robust across demographic segments
rather than an artifact of a particular user profile.

Each injected event is designed to produce \emph{correlated multi-domain
signatures}. For example, user\_01's factory layoff at Day~40 should
simultaneously affect: (a) \appname{Dailyn}: sudden income loss, spending
pattern shifts; (b) \appname{Mealens}: dietary degradation (cheaper,
less nutritious meals); (c) \appname{Ururu}: mood decline, anxiety signals;
(d) \appname{Narrus}: shift in reading topics toward job search, financial
advice. A single-domain analysis can detect the within-domain anomaly; only
cross-domain analysis can identify the \emph{causal pattern} connecting all
four signals.

\subsection{Data Generation}

For each user, we generate 90 daily records per app (360 records per user,
3,600 total), with the following characteristics:

\begin{itemize}[leftmargin=*,itemsep=2pt]
  \item \textbf{\appname{Dailyn}}: 3--8 financial transactions per day with
    category, amount, and memo fields. Amounts are calibrated to each user's
    income level with realistic variance.
  \item \textbf{\appname{Mealens}}: 2--4 meal entries per day with estimated
    nutritional content, meal type (breakfast/lunch/dinner/snack), and
    optional photo description.
  \item \textbf{\appname{Ururu}}: 1--3 mood entries per day with a 1--10
    mood score, free-text journal entry, and detected emotional state tags.
  \item \textbf{\appname{Narrus}}: 0--5 reading sessions per day with
    article/book title, topic tags, reading duration, and highlight/annotation
    count.
\end{itemize}

Pre-event and post-event data distributions differ systematically according to
each user's injected event, creating ground-truth temporal drift patterns that
the system should be able to detect.

\subsection{Models and Hardware}

All experiments use the Qwen3.5 model family exclusively, eliminating
architecture variation as a confound. Table~\ref{tab:models} lists the
model--device assignments.

\begin{table}[H]
\centering
\caption{Model and hardware configuration. All models use Q8 GGUF quantization.
Inference engine: llama.cpp on all devices.}
\label{tab:models}
\begin{tabular}{@{}llll@{}}
\toprule
\textbf{Device} & \textbf{Memory} & \textbf{Model} & \textbf{Role} \\
\midrule
M2 Ultra & 192\,GB & Qwen3.5-35B-A3B-Q8 & Panoramic reasoning (Tier 2) \\
M2 Ultra & 192\,GB & Qwen3.5-9B-Q8      & Scale comparison \\
M2 Ultra & 192\,GB & Qwen3.5-2B-Q8      & Scale comparison \\
M2 Ultra & 192\,GB & Qwen3.5-0.8B-Q8    & IoT baseline \\
M1 Max   & 32\,GB  & Qwen3.5-9B-Q8      & Mid-tier inference (Tier 1) \\
M2 Pro   & 32\,GB  & Qwen3.5-9B-Q8      & Mid-tier inference (Tier 1) \\
\bottomrule
\end{tabular}
\end{table}

The Qwen3.5-35B-A3B is a Mixture-of-Experts (MoE) architecture with 35
billion total parameters but only 3 billion activated per token, providing
knowledge capacity comparable to a 35B dense model at inference costs closer
to a 3B model. This architecture directly embodies the \echostream{} thesis:
a ``medium-scale but highly efficient'' local model.

\subsection{Scoring Methodology}
\label{sec:setup:scoring}

Each generated insight is scored on four dimensions, each on a 0--25 scale,
for a total of 0--100:

\begin{enumerate}[leftmargin=*,itemsep=2pt]
  \item \textbf{Relevance} (0--25): Does the insight address something
    meaningful in the user's life?
  \item \textbf{Specificity} (0--25): Does it reference concrete details
    from the user's data, or is it generic advice?
  \item \textbf{Cross-Domain} (0--25): Does the insight connect patterns
    across multiple life domains?
  \item \textbf{Actionability} (0--25): Can the user take a concrete
    action based on this insight?
\end{enumerate}

Scoring is performed by Claude Opus 4.6 \citep{anthropic2025claude} serving as an
LLM-as-Judge. The judge receives the user's full data context and the
generated insight, and produces scores with justifications. This approach
follows the established practice of using strong LLMs as evaluators
\citep{zheng2023judging}.

The key metric is the \textbf{Insight Increment Ratio} (\iir{}):
\begin{equation}
\iir{} = \frac{S_H}{\frac{1}{4}(S_A + S_B + S_C + S_D)}
\label{eq:iir}
\end{equation}
where $S_H$ is the total score under full panoramic integration (Config~H)
and $S_A, S_B, S_C, S_D$ are scores under single-domain configurations.
An $\iir{} > 1.0$ indicates that cross-domain integration adds value beyond
what any individual domain can provide.


% ══════════════════════════════════════════════════════════════════════════════
\section{Results}
\label{sec:results}
% ══════════════════════════════════════════════════════════════════════════════

\subsection{Experiment A: Cross-Domain Ablation Study}
\label{sec:results:ablation}

We evaluate eight data configurations on all 10 users using
Qwen3.5-35B-A3B-Q8. The configurations range from single-domain (A--D)
through two-domain combinations (E--G) to full panoramic integration (H).
Results are summarized in Table~\ref{tab:ablation}.

\begin{table}[H]
\centering
\caption{Cross-domain ablation results. Scores are averages across 10 users
(0--100 scale). Config~H consistently outperforms all partial configurations.}
\label{tab:ablation}
\begin{tabular}{@{}clc@{}}
\toprule
\textbf{Config} & \textbf{Data Sources} & \textbf{Avg.\ Score} \\
\midrule
A & Finance only (\appname{Dailyn})   & 66.3 \\
B & Diet only (\appname{Mealens})     & 65.1 \\
C & Mood only (\appname{Ururu})       & 63.2 \\
D & Reading only (\appname{Narrus})   & 55.4 \\
\midrule
E & Finance $\times$ Diet             & 76.6 \\
F & Finance $\times$ Mood             & 76.3 \\
G & Diet $\times$ Mood                & 74.1 \\
\midrule
H & Full panoramic (all 4)            & 92.6 \\
\bottomrule
\end{tabular}
\end{table}

The single-domain average is $\bar{S}_{single} = \frac{66.3 + 65.1 + 63.2 + 55.4}{4} = 62.5$, yielding:

\begin{equation}
\iir{} = \frac{92.6}{62.5} = 1.48\times
\end{equation}

The \iir{} ranges from 1.44 (user\_05, e-commerce operations) to 1.53
(user\_02, delivery rider with severe traffic accident). Users experiencing
dramatic cascading life crises consistently show higher \iir{} values,
confirming that cross-domain integration is most valuable when multiple
life domains are simultaneously disrupted. Per-user results are shown
in Table~\ref{tab:iir_per_user}.

\begin{table}[H]
\centering
\caption{Per-user \iir{} values. Users with dramatic cascading life crises
(marked $\geq 1.5$) show higher cross-domain value. Range: 1.44--1.53.}
\label{tab:iir_per_user}
\begin{tabular}{@{}clcc@{}}
\toprule
\textbf{User} & \textbf{Profile} & \textbf{Age} & \textbf{\iir{}} \\
\midrule
user\_01 & Factory worker        & 22 & 1.47 \\
user\_02 & Delivery rider        & 28 & 1.53 \\
user\_03 & Primary teacher       & 31 & 1.45 \\
user\_04 & Grad student          & 26 & 1.48 \\
user\_05 & E-commerce ops        & 27 & 1.44 \\
user\_06 & Freelance illustrator & 29 & 1.46 \\
user\_07 & Hospital resident     & 33 & 1.51 \\
user\_08 & Tech P7 engineer      & 32 & 1.52 \\
user\_09 & Business owner        & 45 & 1.51 \\
user\_10 & Fund manager          & 38 & 1.45 \\
\midrule
\multicolumn{3}{@{}l}{\textbf{Average}} & \textbf{1.48} \\
\bottomrule
\end{tabular}
\end{table}

\paragraph{Dimension-level analysis.} The cross-domain scoring dimension
is the primary driver of the \iir{} gain. Under single-domain configurations
(A--D), the cross-domain score averages 2--5 out of 25, which is expected:
a model with access to only finance data cannot produce cross-domain insights.
Under the full panoramic configuration (H), the cross-domain score rises to
21--24 out of 25. The other three dimensions (relevance, specificity,
actionability) show moderate but consistent improvement, as the richer data
context enables more specific and actionable recommendations.

\paragraph{Two-domain combinations.} The two-domain configurations (E--G)
achieve scores in the 74--77 range, intermediate between single-domain
(55--66) and full panoramic (92.6). This confirms that cross-domain value
scales with the number of integrated domains, though the relationship is
super-linear: the jump from two domains to four domains
($92.6 - 75.7 \approx +17$) is proportionally larger than the jump from one
domain to two domains ($75.7 - 62.5 \approx +13$ for a $2\times$ increase
in data sources).

\paragraph{Qualitative insight examples.}
Consider user\_02 (delivery rider, severe traffic accident at Day~46).
With only \appname{Dailyn} (Config~A), the system detects ``anomalous
spending spikes''---an incomplete observation. With only \appname{Ururu}
(Config~C), it detects ``mood has declined severely.'' With full panoramic
data (Config~H), the system identifies a complete \emph{cascading life crisis}:
income drops to zero for 35 consecutive days, calories crash from 1,800 to
200/day (hospital IV), sleep collapses from 6.2h to 1.5h, and reading topics
shift from casual browsing to legal aid applications and traffic accident
compensation law. The model generates: \emph{``This is a synchronized
multi-domain collapse consistent with a severe physical trauma event.
Financial, nutritional, emotional, and cognitive patterns all pivot around
Day~46. Specific recommendations: apply for work-injury compensation via
[referenced article], increase caloric intake to at least 1,200 kcal
during recovery, consider [specific low-cost career transition paths based
on reading interest patterns].''}

The qualitative difference is categorical: single-domain analysis produces
observations; cross-domain analysis produces \emph{causal explanations} and
\emph{targeted interventions} grounded in the user's specific data trail.


\subsection{Experiment B: Model Scale vs.\ Insight Quality}
\label{sec:results:scale}

To understand the relationship between model scale and insight quality, we run
all 10 users under the full panoramic configuration (Config~H) on four
Qwen3.5 checkpoints, all on the M2~Ultra. Results are shown in
Table~\ref{tab:scale}.

\begin{table}[H]
\centering
\caption{Model scale vs.\ insight quality. All runs use Config~H (full
panoramic) on M2~Ultra 192\,GB. Scores are averages $\pm$ standard deviation
across 10 users.}
\label{tab:scale}
\begin{tabular}{@{}lccc@{}}
\toprule
\textbf{Model} & \textbf{Avg.\ Score} & \textbf{Std.\ Dev.} & \textbf{Quality Tier} \\
\midrule
Qwen3.5-0.8B-Q8     & 48.9 & 4.6 & Unusable \\
Qwen3.5-2B-Q8       & 64.1 & 3.2 & Marginally useful \\
Qwen3.5-9B-Q8       & 79.0 & 3.4 & Good \\
Qwen3.5-35B-A3B-Q8  & 84.6 & 3.4 & Excellent \\
\bottomrule
\end{tabular}
\end{table}

The quality curve exhibits clear diminishing returns:
\begin{itemize}[leftmargin=*,itemsep=2pt]
  \item 0.8B $\to$ 2B: $+15.2$ points ($\Delta = 15.2$)
  \item 2B $\to$ 9B: $+14.9$ points ($\Delta = 14.9$)
  \item 9B $\to$ 35B: $+5.6$ points ($\Delta = 5.6$)
\end{itemize}

The 2B$\to$9B transition represents the steepest quality improvement relative
to the model size increase (4.5$\times$ parameters for $+14.9$ points), while
the 9B$\to$35B transition (3.9$\times$ parameters for $+5.6$ points) yields
only 37.6\% of the per-parameter efficiency. This suggests that for personal
data integration tasks, the practical sweet spot is at the 9B scale: it
achieves 93.4\% of the 35B score ($79.0/84.6$) while requiring substantially
less memory and offering higher throughput.

The 0.8B model's poor performance (48.9) is not merely a matter of degree;
it produces qualitatively different outputs---generic platitudes rather than
data-grounded insights. This establishes a lower bound: sub-2B models are
insufficient for meaningful cross-domain personal reasoning on current
architectures.


\subsection{Experiment C: Device Performance Benchmark}
\label{sec:results:benchmark}

We benchmark inference speed across the three-tier hardware stack.
Table~\ref{tab:benchmark} reports tokens per second (TPS) and time to first
token (TTFT) for each device--model combination.

\begin{table}[H]
\centering
\caption{Inference performance across the three-tier device stack. TPS measured
at two context lengths: long ($\sim$4K tokens) and medium ($\sim$2K tokens).
Inference engine: llama.cpp, quantization: Q8 GGUF.}
\label{tab:benchmark}
\begin{tabular}{@{}llccc@{}}
\toprule
\textbf{Device} & \textbf{Model} & \textbf{TPS (long)} & \textbf{TPS (med)} & \textbf{TTFT (s)} \\
\midrule
M2 Ultra 192G & Qwen3.5-0.8B-Q8    & 137.2 & 135.7 & 0.088 \\
M2 Ultra 192G & Qwen3.5-2B-Q8      & 105.1 & 105.2 & 0.139 \\
M2 Ultra 192G & Qwen3.5-35B-A3B-Q8 & 49.9  & 49.3  & 0.365 \\
M2 Ultra 192G & Qwen3.5-9B-Q8      & 41.3  & 41.4  & 0.471 \\
\midrule
M1 Max 32G    & Qwen3.5-9B-Q8      & 21.9  & 22.0  & 1.138 \\
M2 Pro 32G    & Qwen3.5-9B-Q8      & 12.7  & 12.7  & 1.762 \\
\bottomrule
\end{tabular}
\end{table}

Several observations merit discussion:

\paragraph{MoE efficiency.} The Qwen3.5-35B-A3B-Q8 model achieves 49.9 TPS
despite having 35B total parameters, outperforming the dense 9B model (41.3
TPS) on the same hardware. This validates the MoE architecture's suitability
for on-device deployment: 35B knowledge capacity at super-9B inference speed.

\paragraph{Cross-device scaling.} The same model (Qwen3.5-9B-Q8) exhibits a
clear performance hierarchy: M2~Ultra (41.3 TPS) $>$ M1~Max (21.9 TPS) $>$
M2~Pro (12.7 TPS). The M2~Ultra's advantage stems from its 800~GB/s memory
bandwidth and 76-core GPU, compared to M1~Max's 400~GB/s / 32-core GPU and
M2~Pro's 200~GB/s / 19-core GPU. Memory bandwidth, not compute, is the
binding constraint for LLM inference.

\paragraph{Context length insensitivity.} TPS is nearly identical between
long and medium context lengths across all configurations, indicating that
our workload (personal data summaries, typically 2--4K tokens) falls well
within the efficient operating range of these models.

\paragraph{Practical implications.} Even at the lowest tier (M2~Pro, 12.7
TPS), a 200-token insight is generated in under 16 seconds. At Tier~2
(M2~Ultra with 35B-A3B, 49.9 TPS), the same insight takes 4 seconds. Both
are well within acceptable latency for a personal insight system that
generates reports on demand rather than in real-time conversation.


\subsection{Experiment D: Federation Protocol Evaluation}
\label{sec:results:federation}

We evaluate the federation protocol (Section~\ref{sec:system:federation})
in a real multi-device deployment across a local area network.

\paragraph{Topology.} The M2~Ultra serves as the panoramic query coordinator.
The M1~Max hosts data for \appname{Mealens} (diet) and \appname{Ururu}
(mood). The M2~Pro hosts data for \appname{Narrus} (reading) and
\appname{Dailyn} (finance). All three machines communicate over a standard
home WiFi network.

\paragraph{Protocol.} The panoramic node sends summary requests to both data
nodes in parallel, collects JSON summaries, constructs a unified prompt, and
runs inference on Qwen3.5-35B-A3B-Q8.

Results from 5 runs are summarized in Table~\ref{tab:federation}.

\begin{table}[H]
\centering
\caption{Federation protocol performance over 5 runs. Raw data stays on
origin devices; only summaries are transmitted.}
\label{tab:federation}
\begin{tabular}{@{}lr@{}}
\toprule
\textbf{Metric} & \textbf{Value} \\
\midrule
Raw data on devices (total)      & 108,850 bytes \\
Transmitted summaries (total)    & 867 bytes \\
Compression ratio                & $125.5\times$ \\
\midrule
Federation latency (avg)         & 469.5 ms \\
Federation latency (min)         & 103.2 ms \\
Federation latency (max)         & 1,459.5 ms \\
\midrule
LLM inference time (avg)         & 40.3 s \\
End-to-end latency (avg)         & 40.8 s \\
\midrule
Raw data leaked across devices   & \textbf{0 bytes (all runs)} \\
\bottomrule
\end{tabular}
\end{table}

\paragraph{Compression analysis.} The $125.5\times$ compression ratio
reflects the fundamental design of the summary interface: rather than
transmitting 90 days of individual transactions, meal logs, mood entries,
and reading records, each data node computes domain-specific aggregates
(e.g., weekly spending by category, daily mood trend slopes, top reading
topics) and returns a compact JSON summary. The semantic loss from this
compression is deliberately minimal: the panoramic model needs patterns and
trends, not individual data points.

\paragraph{Latency decomposition.} Federation latency (469.5\,ms average)
is $\sim$1.2\% of the total end-to-end latency (40.8\,s), which is dominated
by LLM inference (40.3\,s). This validates the architectural decision to
separate data collection from inference: the network overhead of federated
data collection is negligible. The max-latency outlier (1,459.5\,ms) likely
reflects WiFi contention and is not architecturally significant.

\paragraph{Privacy verification.} Across all 5 runs, the audit trail confirms
zero raw-data leakage. The data nodes' summary endpoints are stateless: they
receive a time-window parameter and return pre-computed aggregates. No
mechanism exists for the panoramic node to request or receive raw records,
as the API simply does not expose them. This is privacy by \emph{architecture},
not by access control.


% ══════════════════════════════════════════════════════════════════════════════
\section{Discussion}
\label{sec:discussion}
% ══════════════════════════════════════════════════════════════════════════════

\subsection{IIR and the 1.5$\times$ Target}

Our initial target for the Insight Increment Ratio was $1.5\times$. The
observed average of $1.48\times$ approaches but does not fully reach this
target in aggregate, though four out of ten users (user\_02, user\_07,
user\_08, user\_09) individually exceed $1.5\times$. We analyze the
structural factors that shape the IIR:

\paragraph{Scoring saturation.} Single-domain configurations already achieve
moderate scores (55.4--66.3 out of 100) because the relevance, specificity,
and actionability dimensions can be partially satisfied without cross-domain
data. A user's spending patterns alone are enough to generate relevant and
somewhat actionable financial advice. The \iir{} formula divides by this
non-trivial baseline, compressing the ratio.

\paragraph{Cross-domain dimension ceiling.} While the cross-domain dimension
is the primary driver of the \iir{} gain (jumping from 2--5 points in
single-domain to 21--24 in panoramic), this single dimension can contribute
at most 25 out of 100 points. The other three dimensions show more modest
improvements from data integration ($\sim$2--5 points each), limiting the
overall score amplification.

\paragraph{Crisis users exceed the target.} Notably, the four users who
exceed $1.5\times$ (user\_02: 1.53, user\_07: 1.51, user\_08: 1.52,
user\_09: 1.51) all experienced severe cascading life crises---a traffic
accident leading to hospitalization, a medical error resulting in
suspension, a layoff triggering divorce, and a bad debt causing cash flow
collapse. These users exhibit the strongest cross-domain signal because
their crises simultaneously perturb all four data dimensions, creating
correlations that single-domain analysis cannot detect. This supports a
key finding: \textbf{the value of cross-domain integration is greatest
precisely when users need it most---during life crises.}

\paragraph{Synthetic data limitations.} Our synthetic data, while designed
with realistic correlations, may underrepresent the richness of real user
data. In particular, the textual content of mood journal entries and reading
annotations---which would provide the richest signals for cross-domain
reasoning in real data---is necessarily simpler in synthetic generation. We
hypothesize that real user data would yield higher \iir{} values.

The $1.48\times$ average \iir{} represents a substantive improvement.
A 48\% increase in insight quality from data integration---with no change
to the model---is practically significant.
Moreover, the qualitative examples (Section~\ref{sec:results:ablation})
demonstrate that the nature of the improvement is categorical:
single-domain analysis produces observations, while cross-domain analysis
produces explanations and interventions.

\subsection{The Diminishing Returns Plateau}

The model-scale experiment (Section~\ref{sec:results:scale}) reveals that
the practical value curve flattens above 9B parameters. The 9B model achieves
93.4\% of the 35B model's score while running at 84\% of its speed on the
same hardware (41.3 vs.\ 49.9 TPS, noting that the 35B MoE is faster due to
sparse activation).

This finding has a concrete practical implication: \textbf{a user with only a
32\,GB device (M1~Max or M2~Pro) can run a 9B model and achieve near-optimal
insight quality.} The 35B MoE model on a 192\,GB machine provides marginal
improvement. The bottleneck for personal AI quality is not model scale---it
is data availability and integration.

This aligns with the broader \echostream{} thesis: in the personal data
regime, data richness matters more than model size.

\subsection{Social Impact: Early Warning and Preventive Intervention}
\label{sec:discussion:social}

Beyond personal convenience, cross-domain data integration enables a class of
interventions with significant social value. Our synthetic user population
includes several cases where the system's cross-domain analysis could function
as an \emph{early warning system} for cascading life crises.

\paragraph{Cascading crisis detection.}
Consider user\_02 (delivery rider, severe traffic accident at Day~46). The
single-domain mood analysis (Config~C) detects ``mood has declined.'' But the
full panoramic analysis (Config~H) detects a synchronized collapse across
all four domains---income drops to zero, calories crash from 1,800 to 200/day,
sleep collapses from 6.2h to 1.5h, and reading topics shift from casual
browsing to legal aid and accident compensation---identifying a \emph{life
crisis cascade} rather than an isolated mood dip. Similarly, user\_09
(business owner) shows the system detecting a debt-induced health crisis
trajectory: financial stress $\to$ insomnia $\to$ cardiac symptoms $\to$
hospitalization, visible only when financial, emotional, dietary, and reading
data are analyzed together.

\paragraph{Pre-crisis intervention window.}
Several users exhibit \emph{foreshadowing signals} in the days before their
crisis events. User\_08 (P7 engineer) shows rising anxiety in mood data and
shifts in reading topics (from technical articles to layoff-related content)
10 days before the actual layoff. User\_07 (hospital resident) shows
accumulating sleep deprivation and stress escalation across 3 consecutive
night shifts before the medication error. A deployed system could detect
these multi-domain warning patterns and prompt early intervention---suggesting
rest, financial planning, or professional support---before the crisis
materializes.

\paragraph{Psychological support and recovery monitoring.}
Post-crisis, the system can track recovery trajectories across domains.
User\_02's data shows a slow but measurable recovery arc: reading topics
gradually shift from legal aid to career transition, calories slowly
increase from hospital levels, and mood stabilizes at a new (lower)
baseline. This multi-dimensional recovery tracking could complement
professional mental health support, providing therapists with objective
behavioral data rather than relying solely on self-reporting.

\paragraph{Societal implications.}
The ability to detect cascading life crises through passive data collection
has implications for public health infrastructure. In populations where
mental health services are scarce or stigmatized, an on-device system that
can detect distress patterns without requiring the user to actively seek
help---and without exposing their data to any external party---could serve
as a crucial safety net. The privacy-preserving architecture is essential
here: users are more likely to maintain honest records (mood journals,
spending logs) when they trust that no one else will see the data.

\subsection{Practical Deployment Considerations}

\paragraph{Hardware accessibility.} The entire experimental stack uses
consumer Apple hardware. The M2~Pro (the least powerful Tier~1 device) is
available in MacBook Pro configurations starting at approximately \$2,000~USD.
The M2~Ultra Mac Studio (Tier~2) costs approximately \$4,000~USD. While not
inexpensive, these are consumer-grade devices that many knowledge workers
already own or could justify purchasing.

\paragraph{Latency budget.} The end-to-end panoramic query latency of
40.8\,s may seem high for a conversational assistant, but it is appropriate
for the intended use case: periodic (daily or weekly) personal insight
reports, not real-time chat. A user would request a ``weekly life summary''
and receive a comprehensive cross-domain analysis within a minute---a
fundamentally different interaction pattern from cloud chatbots.

\paragraph{Cold start.} The vertical app design addresses cold start through
utility-driven adoption. Users download \appname{Dailyn} because they need
an accounting app, not because they want to train a personal AI. Data
accumulation is a byproduct of normal usage. After 30--60 days of natural
usage across multiple apps, the system has sufficient data for meaningful
cross-domain insights.

\subsection{Limitations}

\paragraph{Synthetic data.} All results are based on synthetic users. While
we have designed the data generation process to be realistic, synthetic data
inevitably underrepresents the complexity, noise, and idiosyncrasy of real
human behavior. A user study with real participants is the critical next step.

\paragraph{LLM-as-Judge.} Our scoring methodology relies on Claude Opus 4.6
as an evaluator. While LLM-as-Judge has been validated in prior work
\citep{zheng2023judging}, it introduces potential biases---particularly the
risk that the judge systematically favors longer or more structured outputs,
which the panoramic configuration naturally produces. Future work should
complement LLM scoring with human evaluation.

\paragraph{Single model family.} All experiments use Qwen3.5, eliminating
architecture as a variable but limiting generalizability. Different model
families (Llama, Mistral, Gemma) may exhibit different model-scale vs.\
insight-quality curves.

\paragraph{Security model.} Our privacy analysis focuses on data
\emph{leakage} (raw data leaving devices). We do not address adversarial
scenarios such as prompt injection attacks on the federation protocol, or
side-channel attacks on the inference process. A production deployment would
require a more comprehensive threat model.

\paragraph{Quantization.} All models use Q8 quantization. While this
provides a good balance of quality and efficiency, the interaction between
quantization level and personal-data reasoning quality has not been explored.
Aggressive quantization (Q4, Q2) may disproportionately degrade the nuanced
reasoning required for cross-domain insight generation.


% ══════════════════════════════════════════════════════════════════════════════
\section{Conclusion}
\label{sec:conclusion}
% ══════════════════════════════════════════════════════════════════════════════

We have presented \echostream{}, a system that demonstrates how cross-domain
personal data integration on consumer hardware produces emergent insights
unreachable by single-domain analysis. Our ablation study across 10 synthetic
users and 8 data configurations establishes an average Insight Increment Ratio
of $1.48\times$, with the cross-domain dimension driving nearly all of the
gain. The model-scale analysis reveals that a 9B model captures 93.4\% of the
35B model's quality, placing the practical sweet spot well within the reach
of consumer 32\,GB devices. The federation protocol achieves $125.5\times$
data compression with zero raw-data leakage, proving that multi-device
personal AI can operate without any data leaving the user's premises.

The broader implication is a paradigm inversion. The dominant cloud AI model
of \emph{large model + small data} reaches an inherent ceiling because users
rationally withhold their most valuable data from remote servers. The
\echostream{} paradigm of \emph{medium model + rich data} exploits an
orthogonal axis of improvement: not larger parameters, but deeper and more
authentic data. As on-device inference capabilities continue to improve with
each hardware generation, this paradigm becomes increasingly viable.

The critical next step is a longitudinal user study with real participants
using real vertical applications. The synthetic data results presented here
are a necessary proof of concept; the true test of the \echostream{} thesis
requires real users generating real data over real time.

\subsection*{Reproducibility}

All code, synthetic data, and experiment scripts are available at
\url{https://github.com/AtomGradient/Prism}. The experimental
results and interactive visualizations are hosted at
\url{https://atomgradient.github.io/Prism}.


% ══════════════════════════════════════════════════════════════════════════════
% REFERENCES
% ══════════════════════════════════════════════════════════════════════════════

\begin{thebibliography}{20}

\bibitem[{Chen et~al.(2025)}]{survey_personalized_llm_2025}
Chen, Z., Wang, Y., Liu, Z., et~al. (2025).
\newblock A survey of personalized large language models.
\newblock \emph{arXiv preprint arXiv:2502.11528}.

\bibitem[{Chuang et~al.(2024)}]{plum2024}
Chuang, Y.-S., Xie, S., Luo, H., Kim, Y., Glass, J., \& He, P. (2024).
\newblock On the way to {LLM} personalization: Learning to remember user
  conversations.
\newblock \emph{arXiv preprint arXiv:2411.13405}.

\bibitem[{Li et~al.(2025)}]{personalllm2025}
Li, T., et~al. (2025).
\newblock {PersonalLLM}: Tailoring {LLMs} to individual preferences.
\newblock In \emph{Proceedings of the International Conference on Learning
  Representations (ICLR 2025)}.

\bibitem[{Wang et~al.(2024)}]{continual_llm_survey_2024}
Wang, T., Zhang, Z., Liang, J., et~al. (2024).
\newblock Continual learning of large language models: A comprehensive survey.
\newblock \emph{arXiv preprint arXiv:2404.16789}.

\bibitem[{Huang et~al.(2024)}]{liber2024}
Huang, Y., Cao, J., Zhang, X., et~al. (2024).
\newblock {LIBER}: Lifelong user behavior modeling based on large language
  models.
\newblock \emph{arXiv preprint arXiv:2411.14713}.

\bibitem[{McMahan et~al.(2017)}]{mcmahan2017fedavg}
McMahan, B., Moore, E., Ramage, D., Hampson, S., \& y~Arcas, B.~A. (2017).
\newblock Communication-efficient learning of deep networks from decentralized
  data.
\newblock In \emph{Proceedings of the 20th International Conference on
  Artificial Intelligence and Statistics (AISTATS)}, pp.\ 1273--1282.

\bibitem[{Ramage \& Mazzocchi(2020)}]{ramage2020federated_analytics}
Ramage, D. \& Mazzocchi, S. (2020).
\newblock Federated analytics: Collaborative data science without data
  collection.
\newblock \emph{Google AI Blog}.

\bibitem[{Hannun et~al.(2023)}]{mlx2023}
Hannun, A., et~al. (2023).
\newblock {MLX}: An efficient machine learning framework for {Apple Silicon}.
\newblock \emph{Apple Machine Learning Research}.

\bibitem[{Gerganov(2023)}]{llamacpp2023}
Gerganov, G. (2023).
\newblock llama.cpp: {LLM} inference in {C/C++}.
\newblock \url{https://github.com/ggerganov/llama.cpp}.

\bibitem[{Anthropic(2025)}]{anthropic2025claude}
Anthropic (2025).
\newblock Claude {Opus} 4.6 model card.
\newblock \url{https://www.anthropic.com}.

\bibitem[{Zheng et~al.(2023)}]{zheng2023judging}
Zheng, L., Chiang, W.-L., Sheng, Y., et~al. (2023).
\newblock Judging {LLM}-as-a-judge with {MT-Bench} and {Chatbot Arena}.
\newblock In \emph{Advances in Neural Information Processing Systems (NeurIPS
  2023)}.

\bibitem[{Qwen Team(2025)}]{qwen2025qwen35}
Qwen Team (2025).
\newblock Qwen3.5 technical report.
\newblock \emph{Alibaba Cloud}.

\end{thebibliography}

\end{document}
